\chapter{From Planck to Dirac}\label{ch:quantum}

\begin{etymology}
\textbf{Quantum} is Latin for ``how much''. Planck used it for the smallest
amount of energy an oscillator can emit or absorb: a packet that cannot be
divided.

\textbf{Spectrum} is Latin for ``appearance'' or ``image'', from \emph{specere},
``to look''. Newton used it for the band of colors a prism makes from white
light. In quantum mechanics, the spectrum of an operator is the set of values a
measurement can return: what the system can show when you look.

\textbf{Uncertainty} comes from \emph{incertus}: \emph{in-}, ``not'', and
\emph{certus}, ``settled'', from \emph{cernere}, ``to separate''. Something that
cannot be sharply separated from what surrounds it.
\end{etymology}

Nobody set out to invent quantum mechanics. Classical physics was forced to give
way by experiments it could not explain, and the fixes proposed between 1900 and
1927 converged on a formalism that Dirac and von Neumann put in its final form.
That formalism is the native language of the qubit.

\section{Three experiments classical physics could not explain}

\begin{history}{1900 \textperiodcentered{} Berlin}
A hot body glows, and the color of the glow depends only on its temperature.
Classical physics predicted that a perfect absorber, a \emph{black body}, would
radiate more and more energy at higher frequencies, without limit; Paul
Ehrenfest later named this the \emph{ultraviolet catastrophe}. Experiments
showed the opposite: the intensity peaks and falls off. On 14 December 1900, Max
Planck presented a formula that fitted the data perfectly. To derive it, he had to
assume that energy is emitted in packets of size $E = h\nu$, proportional to the
frequency. In a letter of 1931 he called it ``an act of desperation''. He did not
believe the packets were real.
\end{history}

\begin{history}{1905 \textperiodcentered{} Bern}
Light shining on a metal knocks electrons out of it, but only if its frequency is
above a threshold; below it, no intensity is enough. Philipp Lenard had measured
this in 1902, and classical waves could not explain it. Albert Einstein, then a
patent clerk, proposed that light itself comes in packets of energy $h\nu$, later
called photons. One photon either has enough energy to free an electron or it
does not. Robert Millikan spent ten years trying to prove Einstein wrong and ended
up confirming him in 1916. Einstein's Nobel Prize of 1921 was for this, not for
relativity.
\end{history}

But Thomas Young had shown in 1801 that light makes interference patterns, like a
wave. Light was behaving as a wave in some experiments and as a particle in others.

\begin{history}{1913 \textperiodcentered{} Manchester and Copenhagen}
Atoms heated in a flame emit light only at specific frequencies, a barcode of
sharp lines unique to each element. Niels Bohr proposed that the electron in a
hydrogen atom can only occupy certain orbits, and emits a photon when it jumps
from one to another, with energy equal to the difference. The model reproduced
the hydrogen lines exactly, and nobody, Bohr included, could say why the orbits
were allowed.
\end{history}

The three cases point the same way: energy in the microscopic world comes in
discrete steps. A theory was needed in which discreteness is built in.

\section{Waves of matter}

In 1924 Louis de Broglie argued in his doctoral thesis that if light, a wave,
behaves like particles, then particles might behave like waves, with wavelength
$\lambda = h/p$. Clinton Davisson and Lester Germer saw electrons diffract off a
nickel crystal in 1927.

Two versions of the new mechanics appeared within months of each other. In 1925
Werner Heisenberg, recovering from hay fever on the island of Helgoland, built a
mechanics out of tables of numbers that Max Born recognized as matrices. In 1926
Erwin Schrödinger wrote a wave equation for a function $\psi(x, t)$:
\[
  i\hbar\,\frac{\partial \psi}{\partial t} = \hat H\psi
  = \Big(-\frac{\hbar^2}{2m}\nabla^2 + V(x)\Big)\psi .
\]
The two looked nothing alike. Schrödinger himself showed in 1926 that they give
the same predictions.

Put a particle in a box of width $L$ with impenetrable walls. The solutions of
Schrödinger's equation are sine waves that fit exactly in the box, and only some
energies are allowed:
\[
  \psi_n(x) = \sqrt{2/L}\,\sin(n\pi x / L), \qquad
  E_n = \frac{n^2 \pi^2 \hbar^2}{2 m L^2}, \qquad n = 1, 2, 3, \dots
\]
Discreteness is no longer a postulate. It comes out of the equation, from the
requirement that the wave fit.

\begin{history}{1926 \textperiodcentered{} Göttingen}
What is $\psi$? Schrödinger thought it was a real wave of matter. Max Born, in a
paper on collisions, proposed instead that it gives probabilities. The famous
rule (the probability is the square of the amplitude) appears in a footnote
added while the paper was in proof. The main text had said the probability was
proportional to $\psi$ itself. Born got the Nobel Prize for it in 1954.
\end{history}

\section{Dirac and von Neumann}

Paul Dirac's \emph{Principles of Quantum Mechanics} (1930) unified matrices and
waves as two descriptions of the same abstract objects: states and operators.
John von Neumann's \emph{Mathematical Foundations of Quantum Mechanics} (1932)
gave it the rigorous form still used today, built on Hilbert spaces
(Chapter~\ref{ch:hilbert}). The bra--ket notation, $\ket{\psi}$ for a state and
$\bra{\phi}$ for its dual, came later, in a short paper Dirac published in 1939.

In that language, quantum mechanics fits in five postulates:

\begin{enumerate}
\item The state of a system is a unit vector $\ket{\psi}$ in a Hilbert space.
\item Every measurable quantity is a Hermitian operator.
\item The possible results of a measurement are the eigenvalues of that operator.
\item After a measurement that returns the eigenvalue $\lambda_k$, the state
becomes the corresponding eigenvector $\ket{\lambda_k}$.
\item Between measurements, the state evolves as
$\ket{\psi(t)} = e^{-i\hat H t/\hbar}\ket{\psi(0)}$.
\end{enumerate}

Every quantum algorithm is these postulates at work. Preparing $\ket{0}$ on every
qubit is preparing a known state. Applying gates is applying
$e^{-i\hat Ht}$ for carefully chosen Hamiltonians. Reading the result is a
measurement. Chapter~\ref{ch:state} goes through them again with one qubit.

\section{The uncertainty principle}

\begin{history}{1927 \textperiodcentered{} Copenhagen}
Heisenberg's first explanation of uncertainty was a thought experiment: to see
where an electron is, you hit it with light, and the light kicks it, disturbing
its momentum. The same year, Earle Kennard proved the precise inequality
$\Delta x\,\Delta p \geq \hbar/2$, where $\Delta$ is the standard deviation over
many identically prepared systems. In 1929 Howard Percy Robertson generalized it to
any pair of quantities.
\end{history}

Heisenberg's microscope is the wrong explanation, and it is still the one most
people learn. The disturbance story suggests that the particle has a definite
position and momentum and that measuring spoils one of them. It does not have
them. The real reason is that position and momentum are operators that do not
commute:
\[
  [\hat x, \hat p] = i\hbar .
\]
Two operators that do not commute have no common set of eigenvectors, so there is
no state in which both have a definite value. For any two quantities $A$ and $B$,
Robertson's inequality says
\[
  \Delta A\,\Delta B \geq \tfrac12 \big|\langle [\hat A, \hat B] \rangle\big| .
\]
A Gaussian wave packet reaches the minimum exactly: squeeze it in position and it
spreads in momentum, and the product stays at $\hbar/2$. This is a statement
about the geometry of the state space, not about instruments.

\section{What this means for a qubit}

\textbf{Measurement destroys superposition.} Measuring forces the system into an
eigenstate of what you measured. You gain a definite value and lose the
superposition you had.

\textbf{Superposition is not ignorance.} Before a measurement, the qubit is not in
a definite state that we happen not to know. In 1964 John Bell turned this into a
testable inequality; Alain Aspect's experiments of 1982 and many since have
confirmed that no theory of hidden definite values reproduces quantum
mechanics. Aspect, John Clauser and Anton Zeilinger shared the 2022 Nobel Prize
for these tests. This is what separates a quantum computer from a classical
machine flipping coins.

\textbf{Unknown states cannot be copied.}

\begin{history}{1982 \textperiodcentered{} FLASH}
Nick Herbert proposed a device he called FLASH that would use quantum
correlations to send signals faster than light. The scheme needed to make many
copies of an unknown quantum state. Refuting it, William Wootters and Wojciech
Zurek, and independently Dennis Dieks, proved in 1982 that no physical process
can copy an arbitrary unknown state.
\end{history}

\begin{concept}{No-cloning theorem}
There is no operation that takes an unknown state $\ket{\psi}$ and a blank qubit
and produces $\ket{\psi}\ket{\psi}$. The proof uses only linearity: a machine that
copied $\ket{0}$ and $\ket{1}$ correctly would turn their superposition into an
entangled state, not into two copies.

\textbf{What it fixes:} it is why quantum key distribution is secure, and why
Herbert's machine cannot work.

\textbf{What it costs:} every classical way of protecting data against errors
involves copies. Quantum error correction had to be invented from scratch
(Part~\ref{part:qec}).
\end{concept}

No-cloning is often presented as a consequence of the uncertainty principle. It
is better to see it as a consequence of linearity. The two are consistent (if you
could clone, you could measure position on one copy and momentum on another), but
the proof needs nothing more than the fact that quantum evolution is linear.
